\section{Điều khảo sát chưa giải quyết}
\label{sec:ch10-dieu-chua-giai-quyet}

Mục trước khép lại bằng một đề xuất mà sách đọc là chưa đủ, nên chỗ đứng của
chương này với khảo sát cần được ghi lại rõ trước khi sang chương sau.

Cái khảo sát đưa được là một tấm bản đồ có số. Nó cho biết có 48 kỹ thuật, cho
biết chúng dồn vào nhóm vấn đề nào, cho biết chúng can thiệp vào bước nào của
quy trình, và quan trọng hơn cả, nó ghi lại bằng lời của chính người trong nghề
rằng không có chuẩn để so chúng với nhau, rằng một nửa không có mã nguồn, và
rằng thước đo hay được chọn sao cho khẳng định được đóng góp đã tuyên bố. Không
một điều nào trong số đó cần một thí nghiệm mới; chúng là quan sát về cách một
lĩnh vực làm việc, và khảo sát ở vị trí tốt để nêu chúng.

Cái khảo sát không đưa được thì gom lại thành ba điều. Thứ nhất, nó không đo gì,
nên mọi phát biểu ở chương này về việc LIME không ổn định đều là phát biểu được
quy qua nó về những bài báo nó dẫn. Thứ hai, nó không xếp hạng, nên câu hỏi
trong tựa đề vẫn không có đáp án sau 18 trang. Thứ ba, bảng tính chất đánh giá
của nó không kèm thủ tục tính, nên nó không phải chỗ để đi tìm một dụng cụ đo.

Ngoài ba điều ấy còn một câu hỏi mà khảo sát không đặt ra, nên chương này cũng
không thể lấy câu trả lời từ nó. Cả 48 kỹ thuật đều là biến thể của một phương
pháp, và toàn bộ phần
thảo luận nói về nhánh LIME như thể vấn đề là vấn đề của nhánh LIME: chuẩn thực
hành cho nghiên cứu về LIME chưa được thiết lập, một khung đánh giá riêng cho
LIME thì nên được xây. Câu hỏi mà khảo sát không hỏi là liệu cùng hình dạng ấy
có lặp lại ở những nhánh khác hay không.

Có lý do để nghi rằng có. Ba quan sát của mục~\ref{sec:ch10-khong-co-cach-so-sanh}
không có chỗ nào riêng cho LIME: so bản sửa của mình với đúng bản gốc, chọn thước
đo xác nhận được đóng góp, và không công bố mã nguồn là ba thói quen mà bất kỳ
nhánh nào cũng có thể mắc. Chính khảo sát đã hé điều đó ở một câu, khi gọi vấn
đề chọn thước đo là phổ biến trong cả cộng đồng XAI chứ không riêng nhánh
LIME~\autocite{p22whichlime}. Nhưng một câu hé ra không phải một bằng chứng, và
chương này dừng ở chỗ nêu nghi vấn.

Chương~\ref{ch:tro-choi-attribution} nhận đúng nghi vấn ấy, và nhận nó với một
bài báo đặt câu hỏi ở mức cả lĩnh vực chứ không ở mức một phương pháp. Nếu hình
dạng của vòng trong hình~\ref{fig:ch10-vong-bien-the} lặp lại ở đó, thì cái sách
đang đọc không còn là một khiếm khuyết của LIME mà là cấu trúc khuyến khích của
việc nghiên cứu attribution.

\index{LIME!biến thể|)}%
